\section{具体行政行为概述}
\begin{enumerate}
  \item 理解性。必考。
  \item 叫书一般就是官方参与了，调解\zd{书}、判决书；起诉\zd{状}。
  \item \zd{具体行政行为 \yjt 直接起诉}。\zd{出题表里不一}，不要看名字看内容。
  \item 四个特点：
        \begin{itemize}
          \item 特定性：\zd{针对特定对象}。
          \item 处分性：\zd{必然影响（主客观相统一）权利义务}。
          \item 外部性：内外有别（\zd{看身份}）。
          \item 行政性：管理性。
        \end{itemize}
  \item 区分：
        \begin{enumerate}
          \item 事实行为：\zd{打人}、行政指导、调解（温柔的建议）、重复处理、过程性行为（告知书、公告、\zd{扣分}）。
          \item 司法行为。司法仲裁、刑事司法行为、行政协助司法行为。
          \item 国家行为、政治行为。
          \item 公安局的区分：客观授权是刑诉，主观目的是刑事。
        \end{enumerate}
  \item 分类：
        \begin{enumerate}
          \item \zd{依职权，依申请}；
          \item \zd{羁束、裁量}：无/有选择余地；合法性/合法性\&合理性评价；不能/可以调解。
          \item \zd{受益、负担}：给好处；给我义务，让我不爽的。
          \item \zd{要式、不要式}。书面盖章；口头，口头传唤。
        \end{enumerate}
  \item 成立（事实判断）和效力（价值判断，有效无效，产生应有的法律效果）。\zd{成立之时是起诉起算}。
        \begin{enumerate}
          \item 作为收到、送达。不作为：
          \item 积极的不作为：明确拒绝。收到文书。
          \item 消极的不作为：不理不睬。默示拒绝。一般等 2 月，紧急情况当场。
        \end{enumerate}
  \item 效力。拘束力（民、官、其他官得遵守）、确定力/不可争力（\zd{确定后不轻易更改}）、执行力（题目带执行二字）。
        \begin{kk}{违法：主观题必占 1 问}
          各自独立。
          \begin{enumerate}
            \item 主体违法：纵向越权、横向越权、地域越权。
            \item 事实依据违法。钓鱼执法。欺骗方式获取证据。
            \item 法律依据违法。
            \item 程序违法。
            \item 明显不合理：\zd{显失公正、滥用职权、畸轻畸重}。拟制为违法。
          \end{enumerate}
        \end{kk}
  \item 效力终止：无效、可撤销、废止、其他原因。
        \begin{enumerate}
          \item 无效：自始当然绝对。\zd{利害关系人}任何时候有权机关（自己，上级）申请宣告无效。
                \zd{2015.5.1 前做出的行为提起诉讼要求宣告无效的，不能受理不能立案}。
          \item 可撤销：\zd{撤销前有效}，撤销时溯及既往。要\zd{赔偿}。
          \item 无效、可撤销区别：违法程度不同。无效：\zd{无（行政主体）资格、无（法律）依据、无可能（去杀人）}。
          \item 废止：法律法规调整、目的实现、公共利益等，\zd{合法而终止}；废止之日失去效力；要\zd{补偿}。
          \item 其他原因：人死了、植物人等。
        \end{enumerate}
\end{enumerate}
